% =====================================================================
%  Section 2 -- Growth of the recurrence
%  CT-S39-DRAFT-A.  Draft only.
%
%  Every citation in 2.0 was verified against the held page image or PDF
%  named beside it, and NOT against OUTLINE.md or a claim text
%  (CT-S39-DRAFT-A/scripts/check_group_a.py; PREREG.md O1-O3).
%  Symbols follow section 1: psi_n, psi_infty, nu_i.
% =====================================================================

\section{Growth of the recurrence}\label{sec:growth}

\subsection{Attribution}\label{ssec:attribution}

The growth statement of this section is, in its leading term, classical. This
subsection says precisely which part is classical, to whom it is due, and what
is left. It is placed first because the section would otherwise read as a claim
of novelty that it does not make.

\paragraph{The radius is prior art.} Write \eqref{eq:n-rec} in the homogeneous
form $Q_{n+1}+a_nQ_n+b_nQ_{n-1}=0$ with $a_n=-b(n+1)\sim-\beta_d n^{d}$ and
$b_n=-1$, so that in the notation of the Newton--Puiseux diagram for such a
recurrence $\alpha=d$ and $\beta=0$. The vertex $P_1=(1,d)$ then lies strictly
above the segment $P_0P_2$, which is case~(a) of Gautschi's Theorem~2.3,
attributed on that page to Perron and Kreuser. Case~(a) supplies two linearly
independent solutions whose consecutive ratios behave like $-a_n$ and
$-b_n/a_n$; for \eqref{eq:n-rec} the first of these gives
$Q_{n+1}/Q_n\sim\beta_d n^{d}$, and the radius statement of
Theorem~\ref{thm:radius} below is an immediate corollary. The minimal/dominant
dichotomy that makes hypothesis (A3) meaningful is Pincherle's, in Gautschi's
Theorem~1.1 on p.~31: the associated continued fraction converges if and only if
the recurrence possesses a minimal solution with non-vanishing initial term.

The attribution is to Gautschi, \emph{SIAM Review} \textbf{9} (1967),
Theorem~2.3(a) p.~35 for Perron--Kreuser and Theorem~1.1 p.~31 for Pincherle;
the same dichotomy appears at DLMF~\S3.6(ii). That the recurrence of this paper
falls in case~(a) is \texttt{CT-S8-LITCHECK-002} (PROVEN); the readings of the
two Gautschi pages are \texttt{CT-S8-LITCHECK-001} and
\texttt{CT-S8-LITCHECK-008} (both INSPECTED).

\paragraph{DLMF (2.9.5)--(2.9.6), and why they do not apply.} It would be
natural to look for the exponent in DLMF~\S2.9(i), which for
$w(n+2)+f(n)w(n+1)+g(n)w(n)=0$ gives the characteristic roots $\rho_j$ by
(2.9.5) and then the exponents explicitly by
\[
  \alpha_j=\frac{f_1\rho_j+g_1}{f_0\rho_j+2g_0}
  \tag{DLMF 2.9.6}
\]
That formula is available \emph{only} under the bounded-coefficient hypothesis
$f(n)\sim\sum_s f_s n^{-s}$, $g(n)\sim\sum_s g_s n^{-s}$ with $g_0\neq0$. Here
$f(n)=-b(n+2)$ grows like $n^{d}$, so for every $d\ge1$ that expansion does not
exist, $f_0$ is not defined, and the clauses do not apply. This is not an
oversight in the source: DLMF routes exactly this class elsewhere, at (2.9.14),
``in which $P$ and $Q$ are any integers see Wong and Li (1992a).'' Both equation
numbers and the deferral are present in the held capture
\texttt{CT-S8-LITCHECK/refs/dlmf\_2\_9.txt}
(\texttt{CT-S8-LITCHECK-003}, INSPECTED).

\paragraph{Wong--Li II was not obtained.} The article DLMF defers to ---
J.~Wong and R.~Li, \emph{Asymptotic expansions for second-order linear
difference equations, II}, \emph{Stud.\ Appl.\ Math.} \textbf{87} (1992)
289--324, DOI \texttt{10.1002/sapm1992874289} --- is paywalled and returned HTTP
403; it has not been read, and nothing in this paper is attributed to it
(\texttt{CT-S8-LITCHECK-004}, INSPECTED). The asymmetry should be stated with
the admission. The unread article cannot reach the exponent $s$, which is
derived in \S\ref{ssec:growth-law} and confirmed independently in
\S\ref{ssec:numerics-growth}, and it cannot reach the framing of
\S\ref{sec:local}--\S\ref{sec:leaf}; it could bear only on whether a closed form
for the amplitude $K$ is already known.

\paragraph{The connection constant.} Classical difference-equation asymptotics
of this kind determines a solution only up to a multiplicative connection
coefficient. That is exactly the gap the amplitude $K$ of
\eqref{eq:growth-law} fills, with the single transcendental input isolated into
$\psi_\infty$ and hypothesis (A3). Two independent sources state the limitation.
Odlyzko, \emph{Asymptotic Enumeration Methods}, p.~77, writes that
\begin{quote}
Wimp and Zeilberger point out that it is usually easy to show that the largest
of the terms of the form (9.35) does show up with a nonzero coefficient, and so
determines the asymptotics of $a_n$ up to a multiplicative constant. However,
the Birkhoff--Trjitzinsky method does not in general provide any techniques for
determining that constant.
\end{quote}
The attribution inside this passage divides, and the division matters: the first
sentence is Odlyzko \emph{reporting Wimp and Zeilberger}, his reference [384],
and only the second is his own assertion. Citing the whole passage to Odlyzko
would misattribute the first half. Independently, Nemes (2026),
arXiv:2508.15246v2, describes connection coefficients as ``often treated as
unknowns in other methods'' and otherwise computed numerically. The two are kept
side by side rather than one replacing the other
(\texttt{S24-12}, \texttt{S24-13}, \texttt{CT-S8-LITCHECK-010}, all INSPECTED).

\paragraph{Nothing here rests on Birkhoff--Trjitzinsky.} The same page of
Odlyzko records the standard objection, that the Birkhoff--Trjitzinsky results
``may not be rigorous, since gaps are alleged to exist in the complicated
proofs'' [211, 383]. A reader who checks p.~77 for the quotation above will read
this too, so it is answered here rather than later. The objection does not bear
on this paper. The radius is Perron--Kreuser through Gautschi; the exponent $s$
is derived directly in \S\ref{ssec:growth-law}; the upper and lower bounds on
$\psi_\infty$ are proved from the recurrence itself. Birkhoff and Trjitzinsky,
\emph{Acta Math.} \textbf{60} (1933), Fundamental Theorem \S9 p.~69, are cited
in this paper for one thing only --- the provenance of the \emph{shape}
$\rho^{n}n^{s}(1+O(1/n))$ of a normal solution --- and that citation is not
load-bearing, because the shape is established here for \eqref{eq:n-rec} by
direct computation (\texttt{CT-S8-LITCHECK-005}, INSPECTED).

\paragraph{What is claimed.} Narrowly: the exponent
$s=\beta_{d-1}/\beta_d+(d-1)/2$ written at the level of the coefficients of $b$
for this recurrence, and the amplitude $K$ of \eqref{eq:K} made explicit through
$\psi_\infty$ and the roots of $b$. The Borel-plane reading of the radius in
\S\ref{ssec:radius} is an expository repackaging of the classical statement and
is presented as one.

\subsection{Convergence of the normalised sequence}\label{ssec:convergence}

\begin{proposition}[\texttt{S1-01}, PROVEN]\label{prop:exists}
Assume (A1). Then $\psi_n=Q_n/B_n$ converges, with no hypothesis on the limit,
and $\psi_\infty-\psi_n=O(n^{-(2d-1)})$.
\end{proposition}

The proof is elementary and is the reason the method statement of this paper
does not invoke a minimal-solution argument. Dividing \eqref{eq:n-rec} by $B_n$
gives the exact recurrence
\begin{equation}\label{eq:psi-rec}
  \psi_n=\psi_{n-1}+\frac{\psi_{n-2}}{b(n-1)b(n)}\qquad(n\ge2),
\end{equation}
whose increments are absolutely summable because $\deg b\ge1$ makes
$\sum_n 1/|b(n-1)b(n)|$ convergent; the telescoping series then converges
absolutely. Nothing is assumed about the limit at this stage.

\begin{remark}[\texttt{S1-05}, STRUCTURAL]\label{rem:pincherle}
Hypothesis (A3) says exactly that $Q_n$ is not the minimal solution of
\eqref{eq:n-rec}, and the minimal subspace has codimension one, so (A3) is
generic. This is Pincherle's dichotomy as recalled in \S\ref{ssec:attribution}.
The remark is interpretive: it is recorded because it explains what (A3) means,
and nothing below depends on it.
\end{remark}

\subsection{Two-sided bounds under positivity}\label{ssec:bounds}

\begin{proposition}[\texttt{S1-02}, PROVEN]\label{prop:positive}
Assume (A1) and in addition $b(n)>0$ for every integer $n\ge1$. Then $\psi_n>0$,
the sequence $(\psi_n)_{n\ge1}$ is strictly increasing, and
$\psi_\infty\ge\psi_1=1$. In particular (A3) holds automatically.
\end{proposition}

\begin{proposition}[Proposition U; \texttt{S26-01}, PROVEN]\label{prop:upper}
Under the hypotheses of Proposition~\ref{prop:positive}, for every $n\ge1$
\[
  \psi_n\le\exp\!\Big(\sum_{k=1}^{n-1}\frac{1}{b(k)b(k+1)}\Big),
  \qquad\text{hence}\qquad
  1\le\psi_\infty\le\exp\!\Big(\sum_{k\ge1}\frac{1}{b(k)b(k+1)}\Big)<\infty .
\]
\end{proposition}

Proposition~\ref{prop:upper} is proved from \eqref{eq:psi-rec} in four steps:
monotonicity gives $\psi_{n-2}\le\psi_{n-1}$; substituting into
\eqref{eq:psi-rec} gives $\psi_n\le(1+\varepsilon_n)\psi_{n-1}\le
\psi_{n-1}e^{\varepsilon_n}$ with $\varepsilon_n=1/(b(n-1)b(n))$; iterating down
to $\psi_1=1$ gives the finite-$n$ form; and the summability already used for
Proposition~\ref{prop:exists} permits passage to the limit.

Two points are worth stating explicitly, because both have been got wrong in
this project's own records. First, the upper bound is \emph{not} part of
Proposition~\ref{prop:positive}: that proposition supplies the monotonicity step
and the lower bound only, and the upper half carries \texttt{S26-01}. Second,
the positivity hypothesis is load-bearing rather than decorative. Of six
sign-changing $b$ satisfying (A1), three violate the inequality outright,
including $b(n)=n-7/2$, for which $\psi_n$ becomes non-positive
(\texttt{S26-04}, VERIFIED, at dps 60 to $n=2000$). No claim is made here about
which sign-changing $b$ do satisfy it.

\subsection{The radius}\label{ssec:radius}

\begin{theorem}[\texttt{S1-03}, PROVEN]\label{thm:radius}
Assume (A1) and (A3). Then $\lim_n|Q_n/(dn)!|^{1/n}=|\beta_d|/d^{d}$, so
$\sum_n Q_n w^{n}/(dn)!$ has radius $d^{d}/|\beta_d|$ in $w$, and radius
$\xi_0=d/|\beta_d|^{1/d}$ in the variable $\zeta=w^{1/d}$.
\end{theorem}

The proof in the source is independent of Gamma asymptotics: the step that
would ordinarily call on Stirling's formula is replaced by a combinatorial
multinomial sandwich. This matters for the attribution of
\S\ref{ssec:attribution}: Theorem~\ref{thm:radius} is a corollary of
Perron--Kreuser, but it is also proved here without importing anything.

\subsection{The growth law and the exponent}\label{ssec:growth-law}

\begin{theorem}[\texttt{S1-04}, PROVEN]\label{thm:growth}
Assume (A2) and (A3). Then
\begin{equation}\label{eq:growth-law}
  \frac{Q_n}{\Gamma(dn+1)}
  =K\Big(\frac{\beta_d}{d^{d}}\Big)^{n}n^{s}\bigl(1+O(1/n)\bigr),
  \qquad
  s=\frac{\beta_{d-1}}{\beta_d}+\frac{d-1}{2},
\end{equation}
with
\begin{equation}\label{eq:K}
  K=\frac{\psi_\infty\,(2\pi)^{(d-1)/2}}{\sqrt{d}\,\prod_{i=1}^{d}\Gamma(1-\nu_i)} .
\end{equation}
\end{theorem}

One classical input is imported, and it is named at the point of use rather than
absorbed silently: the normalisation $h_\infty=1$ for the Gamma ratio, DLMF
(5.11.12), $\Gamma(z+a)/\Gamma(z+b)\sim z^{a-b}$. The decay
$h(n)/h(n-1)=1+O(1/n^{2})$ that reduces the ratio to that single constant is
proved in the source from the functional equation alone. A reader who wishes to
see exactly what is imported should look there and nowhere else.

\begin{remark}[Scope]\label{rem:scope}
Only $\beta_d$ and $\beta_{d-1}$ appear in \eqref{eq:growth-law}. The remaining
coefficients $\beta_0,\dots,\beta_{d-2}$ do not affect the exponent at all; they
enter only through $\psi_\infty$ and the roots $\nu_i$ in \eqref{eq:K}. This is
a statement about the exponent, and it is confirmed rather than assumed in
\S\ref{ssec:numerics-growth}.
\end{remark}

Hypothesis (A2) is consumed here and not before: it is exactly the condition
that no $\nu_i$ is a positive integer, hence that every $\Gamma(1-\nu_i)$ in
\eqref{eq:K} is finite. A statement of \eqref{eq:K} citing only (A1) has omitted
a hypothesis.

\subsection{Numerical confirmation}\label{ssec:numerics-growth}

The exponent was checked independently of its derivation. For all eight quartic
representatives of the $d=4$ family, the measured exponent agrees with
$s=\beta_3/\beta_4+3/2$ to a maximum discrepancy of $7.118\times10^{-16}$, after
four levels of ratio-2 Richardson extrapolation on the ladder
$n=1250,2500,5000,10000,20000$ at dps 80; the pre-registered falsifier
$|\Delta s|>10^{-6}$ did not fire (\texttt{S2-01}, VERIFIED). The accompanying
vacuity arms confirm that the measured exponent does move when $\beta_{d-1}$ is
varied and does not move when the lower coefficients are, which is what
Remark~\ref{rem:scope} asserts.

For the quadratic case $b(n)=3n^{2}+n+1$ the limit itself is available to high
precision:
\[
  \psi_\infty=1.016604441058510936337097213267657686917983514389\ldots
\]
computed to 2317.8 confirmed digits by an $N$-versus-$2N$ agreement criterion at
dps 2800 (\texttt{S3B-02} and \texttt{S3B-09}, both VERIFIED). Hypothesis (A2)
holds for this $b$, its roots being a complex-conjugate pair and in particular
not positive integers (\texttt{S3B-03}, PROVEN). A PSLQ search for an integer
relation between $\psi_\infty$ and a ten-element basis of standard constants
returned NULL at dps 500, 1000 and 2000 with height bound $10^{6}$
(\texttt{S3B-08}, VERIFIED); this is a negative result against a stated search
space and is not evidence that no closed form exists.
